
%--------------------------------------------------
\section{Discusión}

El principal aporte de este trabajo radica en la estimación de la \emph{participación en la billetera} (SoW) y del \emph{tamaño de la billetera} (SioW) en un contexto de información parcial, característico de aplicaciones empresariales reales. En particular, la empresa focal no observa el gasto total de sus clientes en la categoría ni dispone de información directa sobre transacciones realizadas con competidores o de medidas de intención de compra obtenidas mediante encuestas. En este escenario, el modelo \emph{pogit} no se concibe como un mecanismo de predicción directa de variables observables, sino como una herramienta de inferencia estructural que permite reconstruir cantidades latentes coherentes con los patrones transaccionales disponibles.

La aplicación del modelo al segmento de clientes empresariales resulta especialmente pertinente, dado que el entorno B2B presenta dinámicas de compra sustancialmente distintas a las del consumo final. Las compras empresariales suelen caracterizarse por frecuencias relativamente bajas, montos agregados elevados y estructuras de pedido complejas asociadas a descuentos por volumen y relaciones comerciales de largo plazo. En este contexto, la descomposición del gasto observado en un componente de intensidad latente (SioW) y un componente de asignación (SoW) permite distinguir entre el potencial económico total del cliente y el grado en que dicho potencial es capturado por la empresa focal, aportando una lectura económica más rica que la que se obtiene a partir de métricas tradicionales de facturación.

La representación conjunta de las estimaciones ajustadas de SoW y SioW en el plano bidimensional facilita la identificación de perfiles de clientes con implicaciones estratégicas diferenciadas. En particular, se observan clientes con alto SioW y bajo SoW, que representan oportunidades de crecimiento en participación; clientes con bajo SioW y alto SoW, caracterizados por una relación estable pero de alcance económico limitado; y clientes con valores elevados en ambas dimensiones, que constituyen el núcleo estratégico de la cartera. Esta segmentación inducida por el modelo proporciona una base interpretable para la priorización de acciones comerciales y de retención.

Las técnicas de segmentación no supervisada empleadas —k-means y k-medoids— cumplen un rol complementario como herramientas de validación estructural e interpretabilidad. La correspondencia observada entre los clusters definidos a partir de variables transaccionales y las regiones emergentes del espacio SoW--SioW sugiere que las estimaciones del modelo capturan patrones consistentes con agrupaciones naturales de clientes basadas en información observable. Este ejercicio no constituye una validación predictiva en sentido estricto, sino una evaluación de coherencia económica y estabilidad estructural fuera de muestra, adecuada para un modelo cuyos principales objetos de interés son variables latentes no observables directamente.

Finalmente, la comparación entre una especificación del modelo basada en variables transaccionales originales y una alternativa basada en reducción de dimensionalidad mediante PCA pone de manifiesto un trade--off entre parsimonia estadística e interpretabilidad de negocio. Mientras que el uso de componentes principales contribuye a la estabilidad numérica y a la reducción de colinealidad, la especificación con variables originales facilita la traducción directa de los resultados en decisiones operativas. Ambas aproximaciones resultan complementarias y su elección depende del objetivo analítico específico.

%--------------------------------------------------
\section{Trabajos futuros}

Una primera línea de extensión de este trabajo consiste en aplicar la metodología propuesta al segmento de personas naturales. A diferencia de los clientes empresariales, los consumidores finales presentan patrones de compra con mayor frecuencia relativa, menor monto por transacción y una mayor sensibilidad a promociones y estímulos comerciales. Evaluar el desempeño del modelo \emph{pogit} en este contexto permitiría analizar su grado de generalidad y comparar sistemáticamente las distribuciones de SoW y SioW entre segmentos B2B y B2C.

Una segunda extensión metodológica relevante corresponde al desarrollo de modelos que integren simultáneamente clientes empresariales y personas naturales dentro de un mismo marco analítico. Esto podría abordarse mediante modelos de mezcla finita o especificaciones jerárquicas que permitan parámetros diferenciados para los mecanismos de asignación e intensidad según el tipo de cliente. Dichos enfoques permitirían capturar heterogeneidad estructural sin necesidad de estimar modelos completamente separados, manteniendo coherencia estadística y mejorando la eficiencia inferencial.

Otra línea de investigación consiste en incorporar explícitamente la dimensión temporal del comportamiento de compra. La estimación del modelo en ventanas móviles o la inclusión de componentes dinámicos permitiría estudiar la estabilidad temporal de la SoW y del SioW, identificar trayectorias de crecimiento o deterioro en la relación con la empresa focal y diseñar mecanismos de alerta temprana frente a pérdidas potenciales de participación en la billetera.

Desde el punto de vista estadístico, futuros trabajos podrían explorar extensiones distribucionales del componente de intensidad, tales como modelos Poisson--Gamma o binomial negativa, con el fin de capturar posibles fenómenos de sobredispersión no modelados explícitamente en la especificación actual. Estas extensiones permitirían evaluar la robustez de los resultados frente a supuestos distribucionales más flexibles, manteniendo la estructura conceptual del modelo.

Finalmente, las estimaciones de SoW y SioW obtenidas podrían integrarse en marcos de análisis causal o de \emph{uplift modeling} para evaluar el impacto de intervenciones comerciales específicas, como campañas promocionales, cambios en condiciones de descuento o estrategias de fidelización. Esta integración permitiría avanzar desde un análisis descriptivo e inferencial hacia herramientas prescriptivas orientadas a la optimización de decisiones de marketing y a la gestión estratégica del valor del cliente.
